\documentclass[tikz,border=8pt]{standalone}
\usepackage{tikz}
\usepackage{amsmath}
\usepackage{amssymb}
\usepackage{pgfplots}
\pgfplotsset{compat=1.18}
\usepackage{ctex}
\usetikzlibrary{calc,positioning,arrows.meta,matrix,trees}

% LC 54 螺旋遍历
% 4x4 矩阵，红色箭头螺旋路径，4 个边界 top/bottom/left/right

\begin{document}
\begin{tikzpicture}[
  every node/.style={font=\small},
  mcell/.style={draw=gray!60, rectangle, minimum size=0.85cm, thick, fill=white},
  vcell/.style={draw=gray!60, rectangle, minimum size=0.85cm, thick, fill=gray!10},
  ptr/.style={->,>=Stealth,very thick,red!70!black},
  bptr/.style={draw=blue!60,dashed,thick},
]

\def\cs{0.95}

%% 标题
\node[font=\large\bfseries] at (4.0, 1.5)
  {螺旋遍历（LC 54 Spiral Order）};
\node[font=\small,gray] at (4.0, 0.9)
  {红箭头从外到内顺时针，数字为访问顺序；蓝虚线标注四边界};

%% 4x4 矩阵（带访问顺序数字）
%% 顺序: 上左→右(1-4), 右上→下(5-7), 下右→左(8-10), 左下→上(11-12), 内圈同理
%% 矩阵值（= 访问顺序）:
%% [[ 1, 2, 3, 4],
%%  [12,13,14, 5],
%%  [11,16,15, 6],
%%  [10, 9, 8, 7]]

\foreach \r/\c/\v in {
  0/0/1,  0/1/2,  0/2/3,  0/3/4,
  1/0/12, 1/1/13, 1/2/14, 1/3/5,
  2/0/11, 2/1/16, 2/2/15, 2/3/6,
  3/0/10, 3/1/9,  3/2/8,  3/3/7
} {
  \node[mcell,font=\scriptsize\bfseries] at (\c*\cs, -\r*\cs) {\v};
}

%% 边界标注（蓝色虚线矩形，4 层）
% 外圈边界 top=0,bottom=3,left=0,right=3
\draw[bptr] (-0.45, 0.45) rectangle (3.3+0.45-0.05, -3*\cs-0.45);

% 内圈边界 top=1,bottom=2,left=1,right=2
\draw[draw=purple!60,dashed,thick] (0.5, -0.45) rectangle (2.35, -2.4);

%% 边界文字标注（放在图四周空白）
\node[font=\scriptsize,blue!60,anchor=south] at (1.42, 0.52) {top=0};
\node[font=\scriptsize,blue!60,anchor=north] at (1.42, -3*\cs-0.52) {bottom=3};
\node[font=\scriptsize,blue!60,anchor=east,rotate=90] at (-0.52, -1.42) {left=0};
\node[font=\scriptsize,blue!60,anchor=west,rotate=90] at (3.3+0.52, -1.42) {right=3};

\node[font=\scriptsize,purple!60,anchor=south] at (1.42, -0.38) {top=1};
\node[font=\scriptsize,purple!60,anchor=north] at (1.42, -2.47) {bottom=2};

%% 螺旋路径箭头（精确绘制于格中心之间）
% 外圈-上：→ 从(0,0)到(0,3)
\draw[ptr] (0.0+0.42, 0.0) -- (3*\cs-0.42, 0.0);
% 外圈-右：↓ 从(0,3)到(3,3)
\draw[ptr] (3*\cs, 0.0-0.42) -- (3*\cs, -3*\cs+0.42);
% 外圈-下：← 从(3,3)到(3,0)
\draw[ptr] (3*\cs-0.42, -3*\cs) -- (0*\cs+0.42, -3*\cs);
% 外圈-左：↑ 从(3,0)到(1,0)
\draw[ptr] (0.0, -3*\cs+0.42) -- (0.0, -1*\cs-0.42);

% 内圈-上：→ 从(1,1)到(1,2)
\draw[ptr] (1*\cs+0.42, -1*\cs) -- (2*\cs-0.42, -1*\cs);
% 内圈-右：↓ 从(1,2)到(2,2)  -- 注意两格间只有一步
\draw[ptr] (2*\cs, -1*\cs-0.42) -- (2*\cs, -2*\cs+0.42);
% 内圈-下：← 从(2,2)到(2,1)
\draw[ptr] (2*\cs-0.42, -2*\cs) -- (1*\cs+0.42, -2*\cs);
% 内圈-左：↑ 仅1格高无箭头（2,1 → 无继续，已完成）
% 标注最后两格
\node[font=\tiny,red!70!black,anchor=south west] at (1*\cs+0.42, -2*\cs-0.4) {完成};

%% 遍历顺序说明（图右侧）
\node[font=\small\bfseries] at (7.5, 0.2) {遍历顺序};
\node[draw=gray!40,fill=white,thin,rounded corners=2pt,
      font=\scriptsize,inner sep=6pt,align=left]
  at (7.5, -1.5)
  {第1圈（外圈）：\\[2pt]
   上：1, 2, 3, 4\\
   右：5, 6, 7\\
   下：8, 9, 10\\
   左：11, 12\\[4pt]
   第2圈（内圈）：\\[2pt]
   上：13, 14\\
   右：15\\
   下：16};

%% 边界更新规则
\node[draw=gray!40,fill=blue!5,thin,rounded corners=2pt,
      font=\scriptsize,inner sep=6pt,align=left]
  at (7.5, -5.0)
  {\textbf{边界缩进规则：}\\[3pt]
   遍历上边 → top++\\
   遍历右边 → right--\\
   遍历下边 → bottom--\\
   遍历左边 → left++\\[3pt]
   循环条件：top $\le$ bottom \&\& left $\le$ right};

%% 算法说明
\node[draw=gray!50,fill=white,rounded corners=3pt,font=\small,
      inner sep=8pt,align=left]
  at (5.0, -7.2)
  {\textbf{螺旋遍历 O(m×n) 时间，O(1) 空间（输出除外）：}\\[4pt]
   初始化 top=0, bottom=m-1, left=0, right=n-1\\[2pt]
   每轮：按上右下左四方向遍历一圈边界，每遍历一边后缩进对应边界\\[2pt]
   直至 top > bottom 或 left > right 终止};

\end{tikzpicture}
\end{document}
